\documentclass[journal,onecolumn]{IEEEtran}

% ── Encabezado ──────────────────────────────────────────────
\usepackage{fancyhdr}
\pagestyle{fancy}
\fancyhf{}
\renewcommand{\headrulewidth}{0.4pt}
\fancyhead[L]{\small ERS/SRS Final}
\fancyfoot[C]{\thepage}

% ============================================================
% IDIOMA Y CODIFICACIÓN  (mover ANTES de csquotes)
% ============================================================
\usepackage[utf8]{inputenc}
\usepackage[T1]{fontenc}
\usepackage[spanish,es-tabla]{babel}

% ============================================================
% CITAS Y TEXTO
% ============================================================
\usepackage{csquotes}

% ============================================================
% ESPACIADO Y LISTAS
% ============================================================
\usepackage{parskip}
\setlength{\parskip}{4pt}
\usepackage{enumitem}

% ============================================================
% IDIOMA Y CODIFICACIÓN
% ============================================================
\usepackage[spanish,es-tabla]{babel}
\usepackage[utf8]{inputenc}
\usepackage[T1]{fontenc}

% ============================================================
% PAPEL Y GEOMETRÍA
% ============================================================
\usepackage{geometry}
\geometry{a4paper, top=2.5cm, bottom=2.5cm, left=2.5cm, right=2.5cm}

% ============================================================
% IMÁGENES
% ============================================================
\usepackage{graphicx}
\graphicspath{{Figuras/}}

% ============================================================
% TABLAS
% ============================================================
\usepackage{booktabs}
\usepackage{tabularx}
\usepackage{longtable}

% ============================================================
% MATEMÁTICAS
% ============================================================
\usepackage{amsmath}
\usepackage{amssymb}

% ============================================================
% CÓDIGO FUENTE
% ============================================================
\usepackage{listings}

% ============================================================
% FORMATO ACADÉMICO
% ============================================================
\usepackage{microtype}
\usepackage{float}

% ============================================================
% CAPTIONS (TÍTULOS DE FIGURAS/TABLAS)
% ============================================================
\usepackage{caption}
\captionsetup{font=small, labelfont=bf}

% ============================================================
% REFERENCIAS / HIPERVÍNCULOS
% ============================================================
\usepackage[hidelinks]{hyperref}

% ============================================================
% NUMERACIÓN Y NOMBRES EN ESPAÑOL
% ============================================================
\renewcommand{\thetable}{\arabic{table}}
\def\tablename{Tabla}

\renewcommand{\thesection}{\arabic{section}}
\renewcommand{\thesubsection}{\thesection.\arabic{subsection}}
\renewcommand{\thesubsubsection}{\thesubsection.\arabic{subsubsection}}

% ============================================================
% ESTILO DE TÍTULOS DE SECCIÓN
% ============================================================
\usepackage{titlesec}
\titleformat{\section}{\normalfont\Large\bfseries}{\thesection}{1em}{}
\titleformat{\subsection}{\normalfont\large\bfseries}{\thesubsection}{1em}{}
\titleformat{\subsubsection}{\normalfont\normalsize\bfseries}{\thesubsubsection}{1em}{}

% ============================================================
% CADA SECCIÓN PRINCIPAL (\section) INICIA EN PÁGINA NUEVA
% ============================================================
\usepackage{etoolbox}
\pretocmd{\section}{\clearpage}{}{}

% ============================================================
% BIBLIOGRAFÍA (BIBLATEX + BIBER)
% ============================================================
\usepackage[
backend=biber,
style=ieee,
sorting=none
]{biblatex}
\addbibresource{referencias_estructura.bib}

% ============================================================
% DOCUMENTO
% ============================================================
\begin{document}


\section{ERS/SRS Final del Sistema MundiPets}
\label{sec:ers-final}

\subsection{Estructura conforme a ISO/IEC/IEEE 29148:2018}
\label{sec:estructura-29148}

El ERS final de MundiPets se organizó siguiendo la estructura de contenido que
ISO/IEC/IEEE 29148:2018 establece para una Especificación de Requisitos de
Software \cite{ieee29148}. La Tabla~\ref{tab:correspondencia-29148} muestra la
correspondencia explícita entre las cláusulas de la norma y las secciones
efectivamente desarrolladas en el documento del equipo; no se transcribe aquí
el contenido de cada sección, que se encuentra íntegro en el ERS consolidado
del repositorio, sino que se documenta dónde localizarlo y bajo qué cláusula
normativa queda encuadrado.

\begin{table}[H]
	\centering
	\caption{Correspondencia entre las cláusulas de ISO/IEC/IEEE 29148:2018 y las secciones del ERS de MundiPets.}
	\label{tab:correspondencia-29148}
	\footnotesize
	\renewcommand{\arraystretch}{1.3}
	\begin{tabularx}{\textwidth}{
		|>{\raggedright\arraybackslash}p{3.6cm}
		|>{\raggedright\arraybackslash}X|}
		\hline
		\textbf{Cláusula ISO/IEC/IEEE 29148:2018} & \textbf{Sección del ERS de MundiPets} \\
		\hline
		1. Introducción & Sección 1 --- Introducción \\
		\hline
		\quad 1.1 Propósito & 1.1 Propósito del documento \\
		\hline
		\quad 1.2 Alcance & 1.2 Alcance del producto \\
		\hline
		\quad 1.3 Definiciones, acrónimos y abreviaturas & 1.3 Definiciones, acrónimos y abreviaturas (glosario) \\
		\hline
		\quad 1.4 Referencias & 1.4 Referencias \\
		\hline
		\quad 1.5 Visión general del documento & 1.5 Visión general del documento \\
		\hline
		2. Descripción general & Sección 2 --- Descripción general \\
		\hline
		\quad 2.1 Perspectiva del producto & 2.1 Perspectiva del producto y límite del sistema (incluye el entorno operativo) \\
		\hline
		\quad 2.2 Funciones del producto & 2.2 Funciones del producto \\
		\hline
		\quad 2.3 Características de los usuarios & 2.3 Clases y características de usuarios; 2.4 Mapa de stakeholders (matriz poder/interés y conflictos potenciales) \\
		\hline
		\quad 2.4 Restricciones generales & 2.5 Entorno operativo (aspectos no cubiertos en 2.1) \\
		\hline
		\quad 2.5 Suposiciones y dependencias & 2.6 Suposiciones y dependencias \\
		\hline
		3. Requisitos específicos & Sección 3 --- Requisitos específicos completos \\
		\hline
		\quad 3.1 Requisitos de interfaces externas & 3.1 Interfaces externas (usuario, hardware, software) \\
		\hline
		\quad 3.2 Requisitos funcionales & 3.2 Requisitos funcionales (RF-01 a RF-27) \\
		\hline
		\quad 3.3 Requisitos de rendimiento & Incorporados dentro de 3.3 Requisitos no funcionales (RNF-02, RNF-03, RNF-12) \\
		\hline
		\quad 3.4 Restricciones de diseño & 3.5 Restricciones de diseño (RD-01 a RD-09) \\
		\hline
		\quad 3.5 Atributos del sistema software (calidad) & 3.3 Requisitos no funcionales (RNF-01 a RNF-16) \\
		\hline
		\quad 3.6 Otros requisitos (legales, de negocio) & 3.4 Requisitos legales \\
		\hline
		4. Información de apoyo & 3.6 Historias de usuario y criterios de aceptación (verificabilidad); Sección 4 --- Modelado del sistema con UML; Sección 5 --- Priorización y trazabilidad extendida (matriz de trazabilidad) \\
		\hline
	\end{tabularx}
\end{table}

La correspondencia es completa: ninguna cláusula de contenido que la norma
exige para un SRS quedó sin sección equivalente en el ERS de MundiPets. La
distinción entre requisitos de rendimiento (cláusula 3.3) y atributos de
calidad del software (cláusula 3.5) --- que la norma trata como categorías
separadas--- se resolvió integrando ambas dentro de la sección única de
Requisitos no funcionales del documento del equipo, clasificando internamente
cada RNF según corresponda a rendimiento, seguridad, usabilidad, mantenibilidad
u otro atributo de calidad; esta es una decisión de organización del equipo,
no una omisión, y se declara aquí de forma explícita por la misma razón que
exige la verificabilidad de un requisito individual: que la trazabilidad entre
lo declarado y lo exigido sea comprobable por cualquier lector, incluido el
tribunal.

El ERS del equipo contiene, además de las cláusulas normativas anteriores, dos
bloques de contenido que exceden el mínimo que exige ISO/IEC/IEEE 29148:2018:
la Sección 6, Producto Mínimo Viable (alcance del MVP, arquitectura de
despliegue y evidencia de funcionamiento), y la Sección 7, Componente empírico
del proyecto (protocolo de validación experimental, preguntas de investigación
en formato PICOC, hipótesis, variables, participantes, instrumentos, plan de
análisis estadístico y registro previo del protocolo). Ninguno de estos dos
bloques corresponde a una cláusula de la norma y, por lo tanto, no se fuerza su
inclusión en la Tabla~\ref{tab:correspondencia-29148}: son contenido adicional
que el equipo decidió mantener como valor agregado del proyecto integrador,
más allá de la estructura mínima que un SRS conforme a 29148 exige.

\subsection{Requisitos clave y cambios incorporados en PE1--PE5}
\label{sec:cambios-pe1-pe5}

La Tabla~\ref{tab:cambios-clave} documenta la trazabilidad de cambios de los
requisitos del sistema para los que el equipo cuenta con evidencia completa y
verificable del ciclo defecto--corrección: redacción original, defecto
detectado en la inspección de la PE4, corrección aplicada y estado final.
Estos cinco requisitos fueron objeto de una solicitud de cambio (RFC) formal
ante el Comité de Control de Cambios, tramitada durante la PE4 y resuelta de
forma definitiva al cierre de la PE5; dos de ellos --- RF-06 y RNF-02---
corresponden a la categoría \emph{Debe tener} (Must) de la priorización
MoSCoW; los tres restantes --- RF-07, RF-10 y RNF-03--- corresponden a
\emph{Podría tener} (Could) y \emph{Debería tener} (Should)
respectivamente, pero se incluyen porque son los únicos casos, junto a los
dos anteriores, en los que el equipo puede exhibir la redacción exacta antes
y después de la corrección, y no solo declarar que un cambio ocurrió.

\begin{table}[H]
	\centering
	\caption{Requisitos con trazabilidad de cambio completa (redacción original, defecto, corrección) incorporados durante PE1--PE5.}
	\label{tab:cambios-clave}
	\footnotesize
	\renewcommand{\arraystretch}{1.3}
	\begin{tabularx}{\textwidth}{
		|>{\raggedright\arraybackslash}p{1.1cm}
		|>{\raggedright\arraybackslash}X
		|>{\raggedright\arraybackslash}p{2.6cm}|}
		\hline
		\textbf{ID} & \textbf{Cambio aplicado y motivo} & \textbf{Práctica} \\
		\hline
		RF-06
		(Must) &
		\emph{Original:} la salida del requisito era información comparativa
		general para evaluar compatibilidad, junto con una advertencia de
		que el resultado era orientativo. \emph{Defecto (RFC-I, PE4):} el
		propio ERS exigía en RNF-06 y RNF-13 una explicación de la
		recomendación que RF-06 no declaraba producir --- el requisito no era
		verificable contra sus propios atributos de calidad. \emph{Corrección
		(cierre PE5):} la salida se redefinió como un veredicto de
		compatibilidad en tres niveles discretos (compatible, compatible con
		reservas, no recomendado), acompañado de la explicación exigida por
		RNF-13. \emph{Estado final:} corregido y verificado; cubierto en el
		MVP mediante heurística explicable. &
		PE2 (elicitación) $\rightarrow$ PE4 (defecto vía RFC-I) $\rightarrow$
		PE5 (corrección e integración) \\
		\hline
		RF-07
		(Could) &
		\emph{Original:} el requisito permitía la comunicación entre
		propietarios e interesados mediante mensajes, sin mecanismo de
		reporte de abuso. \emph{Defecto (RFC-F, PE4):} el propio ERS
		reconocía el conflicto de ``comunicación inadecuada'' entre usuarios
		sin que ningún requisito funcional lo resolviera. \emph{Corrección
		(cierre PE5):} se incorporó la posibilidad de reportar un mensaje o
		interacción como spam, lenguaje ofensivo o intento de acuerdo externo
		indebido, con registro del motivo asociado a la conversación; esto
		obligó además a redefinir RNF-15 (RFC-H) para que la explicabilidad
		de la moderación fuera coherente con este nuevo mecanismo de reporte
		manual y no con una moderación automática que el sistema no
		implementa. \emph{Estado final:} corregido y verificado. &
		PE2 (elicitación) $\rightarrow$ PE4 (defecto vía RFC-F) $\rightarrow$
		PE5 (corrección e integración) \\
		\hline
		RF-10
		(Should) &
		\emph{Original:} el requisito generaba recordatorios por
		notificación en la aplicación y WhatsApp, sin definir qué canal
		prevalecía si el otro fallaba. \emph{Defecto (RFC-C, PE4):} la
		restricción de diseño RD-08 exigía un canal de respaldo ante caída
		del proveedor externo de mensajería, pero RF-10 no lo formalizaba.
		\emph{Corrección (cierre PE5):} se redefinió WhatsApp como canal
		principal y la notificación in-app como respaldo automático cuando
		el proveedor externo no está disponible, con su correspondiente
		criterio de verificación de conmutación de canal. \emph{Estado
		final:} corregido y verificado; cubierto en el MVP mediante
		simulación del canal de respaldo. &
		PE2 (elicitación) $\rightarrow$ PE4 (defecto vía RFC-C) $\rightarrow$
		PE5 (corrección e integración) \\
		\hline
		RNF-02
		(Must) &
		\emph{Original:} exigía una disponibilidad mínima del 99\,\% mensual
		sin acotar el entorno de medición. \emph{Defecto (RFC-E, PE4):} el
		MVP local del equipo no permite verificar disponibilidad mensual, al
		no operar en un ambiente de uso continuo. \emph{Corrección (cierre
		PE5):} se acotó explícitamente el criterio al entorno de producción
		con backend desplegado, declarando que no es exigible ni medible
		sobre el MVP local. \emph{Estado final:} corregido y verificado. &
		PE3 (especificación) $\rightarrow$ PE4 (defecto vía RFC-E)
		$\rightarrow$ PE5 (corrección e integración) \\
		\hline
		RNF-03
		(Should) &
		\emph{Original:} exigía un tiempo de respuesta máximo de 2 segundos
		en el percentil 95 con 100 usuarios concurrentes, sin acotar el
		entorno de medición. \emph{Defecto (RFC-G, PE4):} el umbral no era
		verificable con la infraestructura disponible para el equipo; la
		propuesta original de relajar el umbral fue rechazada por el CCB de
		cierre de la PE5, que prefirió acotar el alcance en vez de bajar la
		exigencia. \emph{Corrección (cierre PE5):} se acotó el criterio al
		entorno de producción, igual que RNF-02, sin modificar el valor
		objetivo de 2 segundos. \emph{Estado final:} corregido y verificado.
		&
		PE3 (especificación) $\rightarrow$ PE4 (defecto vía RFC-G)
		$\rightarrow$ PE5 (corrección e integración, rechazo de la propuesta
		original) \\
		\hline
	\end{tabularx}
\end{table}

Además de estos cinco, la revisión de calidad de requisitos funcionales
realizada en la Entrega 4 (versión 4.0 del ERS) corrigió por ambigüedad,
duplicidad o inconsistencia otros seis requisitos de la categoría Must ---
RF-02, RF-05, RF-12, RF-17, RF-18 y RF-23---, según consta en el historial de
versiones del documento. El equipo no reconstruye aquí la redacción original
de cada uno de estos seis por no contar con el registro de defecto
individual con el que sí cuenta para los cinco de la Tabla~\ref{tab:cambios-clave}:
declarar un ``antes'' sin ese respaldo sería exactamente el tipo de
evidencia no verificable que la sección de integridad académica de la
rúbrica sanciona con mayor severidad. Lo que sí es verificable, y se declara
como tal, es que esa corrección de calidad ocurrió y está registrada en el
propio ERS como parte del cierre normal de la Entrega 4.

La existencia de esta tabla es en sí misma la evidencia de que el ERS
entregado en la PE5 es la versión final integrada del documento y no una
versión previa a las correcciones: cada fila muestra un requisito que en
algún punto de PE1--PE4 tuvo un defecto real, registrado en un acta de
inspección o en una RFC, y que llegó a la PE5 ya corregido y verificado
contra la versión vigente del ERS (v4.0).

\printbibliography[heading=none]

\end{document}
